\section{Reference Results and Limitations}
\label{sec:benchmarks}
The repository includes \texttt{benchmark\_results.json}, which stores a reference Raspberry Pi 4 profile. The benchmark script uses those values only when emulation is explicitly requested; failed live runs do not generate replacement results. Because the repository does not include a device log, power trace, validation predictions, or calibration images, Table~\ref{tab:reference} is reported as a reference/emulated profile. It must not be interpreted as an independently reproduced physical experiment.

\begin{table}[t]
\centering
\caption{Reference Raspberry Pi 4 profile supplied with the repository.}
\label{tab:reference}
\setlength{\tabcolsep}{4pt}
\begin{tabular}{lcc}
\toprule
\textbf{Metric} & \textbf{FP32 reference} & \textbf{INT8 reference}\\
\midrule
Mean latency (ms) & 220.0 & 95.0\\
90th-percentile latency (ms) & 245.0 & 105.0\\
Throughput (inferences/s) & 4.55 & 10.53\\
Process memory (MB) & 240.0 & 65.0\\
Model size (MB) & 13.6 & 3.3\\
Assumed power (W) & 3.2 & 3.0\\
Derived energy (J/inference) & 0.704 & 0.285\\
\bottomrule
\end{tabular}
\end{table}

Relative to the supplied FP32 reference, the INT8 profile corresponds to a 2.32$\times$ lower mean latency, a 4.12$\times$ smaller model artifact, and a 3.69$\times$ lower reported process-memory figure. The energy values follow directly from the listed latency and assumed power: $0.220\,\mathrm{s}\times3.2\,\mathrm{W}=0.704\,\mathrm{J}$ and $0.095\,\mathrm{s}\times3.0\,\mathrm{W}=0.285\,\mathrm{J}$. They are therefore derived estimates, not power measurements.

The repository does not currently support a defensible claim about ImageNet accuracy: the conversion scripts previously used random calibration arrays, no validation dataset is distributed, and the original accuracy fields were inherited reference values. Similarly, the Jetson Nano and Google Coral rows present in the earlier draft were not accompanied by raw logs or executable platform-specific benchmark code and have been removed from the manuscript. The same applies to claims about batching, concurrent multi-model execution, and container overhead. These are worthwhile future experiments, but they are not established by the current artifact.

This separation improves scientific usability. A reader can exercise conversion and serving after supplying the required dependencies and calibration data, while the missing evidence is visible instead of being hidden behind precise-looking numbers. The repository's benchmark tool should be run in live mode only after the model and calibration data have been supplied; its output should then be archived with the environment and hardware metadata listed in Section~\ref{sec:methodology}.
